\exercisesection{5}

\begin{enumerate}
\def\labelenumi{\arabic{enumi}.}
\item
  Define an endomorphism \(u\) of the \(\mathbf{Z}\) -module \(\mathbf{M} = \mathbf{Z}^{(\mathbf{N})}\) such that, for the prime ideal \(p = (0)\) of \(\mathbf{Z}\) , \(u_{p}\) is an automorphism of \(M_p\) but that there exists no \(f \neq 0\) in \(\mathbf{Z}\) such that \(u_f\) is a surjective endomorphism of \(M_f\) .
\item
  Let K be a field and B the ring of polynomials with coefficients in K in an infinite system of indeterminates \(X_{i}\) ( \(i \in N\) ); let b be the ideal of B generated by the products \(\mathbf{X}_i\mathbf{X}_j\) ( \(i \neq j\) ), A the quotient ring B/b, \(\xi_i\) the class of \(\mathbf{X}_i\) mod. b ( \(i \in \mathbb{N}\) ), p the ideal of A generated by the \(\xi_i\) of index \(i \geqslant 1\) , which is prime, and u the canonical mapping A → A/p.~Show that \(u_p: \mathbf{A}_p \to (\mathbf{A}/\mathfrak{p})_p\) is bijective but that there exists no \(f \in \mathbf{A} - \mathfrak{p}\) such that \(u_f\) is bijective.
\item
  Let A be a ring and P a finitely generated projective A-module. Show that A is isomorphic to a finite product of rings \(\prod_{i\in I}A_{i}\) and P is isomorphic to a product \(\prod_{i\in I}P_{i}\) , where \(P_{i}\) is a projective \(A_{i}\) -module of rank \(n_{i}\) , the \(n_{i}\) being distinct (use Theorem 1 and the fact that \(\operatorname{Spec}(A)\) is quasi-compact).
\item
  Let A be a (commutative) ring and B a (not necessarily commutative) ring containing A. Suppose that the left A-module B is projective and finitely generated. Show that A is a direct factor of B. (Reduce it to the case where A is a local ring and use § 3, no. 2, Corollary 1 to Proposition 5.)
\item
  Let A be a reduced ring and M a finitely generated A-module. Suppose that there exists an integer n \textgreater{} 0 such that, for every homomorphism \(\phi\) of A to a field \(K_{\phi}\) , \([M \otimes_{A} K_{\phi}: K_{\phi}] = n\) . Show that M is a projective module of rank n.~(Reduce it to the case where A is a local ring and use §2, no. 2, Proposition 7.)
\end{enumerate}

¶ 6. Let A be an integral domain, K its field of fractions and M an A-module. Show that the following properties are equivalent:

\((\alpha)\) M is a projective A-module such that \([\mathbf{M} \otimes_{\mathbf{A}} \mathbf{K} : \mathbf{K}]\) is finite.

(β) M is finitely generated projective A-module.

\((\gamma)\) M is finitely generated and, for every maximal ideal m of A, the \(\mathbf{A}_{m}\) -module \(\mathbf{M}_{m}\) is free.

(To show that \((\alpha)\) implies \((\beta)\) , use Algebra, Chapter II, § 5, no. 5, Proposition 9. To show that \((\gamma)\) implies \((\alpha)\) , observe first that M is torsion-free and conclude that, for every prime ideal \(p\) of A, \(M_{p}\) is a free \(A_{p}\) -module of constant rank.)

\begin{enumerate}
\def\labelenumi{\arabic{enumi}.}
\setcounter{enumi}{6}
\tightlist
\item
  Let A be an absolutely flat (commutative) ring (cf.~§ 4, Exercise 16) and let X be its spectrum.
\end{enumerate}

\begin{enumerate}
\def\labelenumi{(\alph{enumi})}
\item
  Let F be a closed subset of X and a the corresponding ideal of A (loc. cit.). Show that M = A/a is a monogenous flat A-module such that \(M_{m}\) is a free \(A_{m}\) -module for all \(m \in X\) , but that M is projective only if F is open.
\item
  Let E be the module the direct sum of the A/m, for \(m \in X\) . Show that E is a flat A-module such that \(E_{m}\) is a free \(A_{m}\) -module of rank 1 for all \(m \in X\) ; for E to be a projective A-module, it is necessary and sufficient that X be finite (note that, if A/m is a projective A-module, m is finitely generated).
\end{enumerate}

\begin{enumerate}
\def\labelenumi{\arabic{enumi}.}
\setcounter{enumi}{7}
\tightlist
\item
  Let A be a ring and M a projective A-module of rank n.~Show that there exists a finitely generated (commutative) A-algebra B such that the A-module
\end{enumerate}

B is faithfully flat and \(\mathbf{M}_{(\mathbf{B})} = \mathbf{M} \otimes_{\mathbf{A}} \mathbf{B}\) is a free B-module of rank \(n\) (use Theorem 1 (d) and Proposition 3). Obtain the converse.

¶ 9. (a) Let A be a ring, \((f_i)_{i \in I}\) a finite family of elements of A generating the ideal A and M an A-module. Suppose given in each of the \(A_{f_i}\) -modules \(M_{f_i}\) an element \(z_i\) such that, for \(i \neq j\) , the canonical images of \(z_i\) and \(z_j\) in \(M_{f_i f_j}\) are equal. Show that there exists a unique \(z \in M\) such that, for all \(i\) , the canonical image of \(z\) in \(M_{f_i}\) is equal to \(z_i\) . (Note that, for every integer \(k > 0\) , there exists a family \((g_i)_{i \in I}\) of elements of A such that \(\sum_{i} g_i f_i^k = 1\) .)

\begin{enumerate}
\def\labelenumi{(\alph{enumi})}
\setcounter{enumi}{1}
\item
  Let A be a ring and P a projective A-module of rank 1. Show that the ring \(\operatorname{End}_{\mathbf{A}}(\mathbf{P})\) is canonically isomorphic to A (use Theorem 3).
\item
  Let A be a ring and P a projective A-module of rank n.~For every endomorphism u of P, there corresponds canonically to \(\wedge^{n}u\) , by (b), an element \(\det(u)\) of A, called the determinant of u, such that \(\det(u \circ v) = (\det u)(\det v)\) for two endomorphisms u, v of P. The element \(\chi_{u}(T) = \det(T.1 - u)\) of A{[}T{]} is called the characteristic polynomial of u (cf.~no. 3, Proposition 4); show that \(\chi_{u}\) is a monic polynomial of degree n and that \(\chi_{u}(u) = 0\) (use (a) and §3, no. 3, Theorem 1). Deduce that, for u to be bijective, it is necessary and sufficient
\end{enumerate}

that \(\det(u)\) be invertible in A. If A is an integral domain, \(\chi_u(T) = \prod_{i=1}^{n} (T - \alpha_i)\) the decomposition into the linear factors of \(\chi_u\) in an algebraically closed extension

of the field of fractions of A, and \(\chi_{q(u)}(T) = \prod_{i=1}^{l} (T - q(\alpha_i))\) for every polynomial \(q \in A[T]\) . Generalize similarly Proposition 14 of Algebra, Chapter VII, § 5, no. 6.

\begin{enumerate}
\def\labelenumi{\arabic{enumi}.}
\setcounter{enumi}{9}
\tightlist
\item
  Let A be a ring and \(E(A)\) the set of classes of finitely generated projective A-modules; let G be the Z-module of formal linear combinations of elements of \(E(A)\) and N the submodule of G consisting of the elements \(\xi - \xi' - \xi''\) , where \(\xi, \xi', \xi''\) are the classes of three projective A-modules M, M', M'' such that there exists an exact sequence
\end{enumerate}

\[
0 \rightarrow \mathrm{M} ^ {\prime} \rightarrow \mathrm{M} \rightarrow \mathrm{M} ^ {\prime \prime} \rightarrow 0
\]

(it amounts to the same to say that M is isomorphic to the direct sum \(M' \oplus M''\) by virtue of Algebra, Chapter II, §2, no. 2). Let \(K(A)\) be the quotient Z-module G/N and, for every finitely generated projective A-module, let \(\kappa(M)\) (or \(\kappa_{A}(M)\) ) be its class in \(K(A)\) .

\begin{enumerate}
\def\labelenumi{(\alph{enumi})}
\item
  Show that there exists on \(K(\mathbf{A})\) a unique commutative ring structure whose addition is that the Z-module \(K(\mathbf{A})\) and multiplication is such that \(\kappa(\mathbf{M})\kappa(\mathbf{N}) = \kappa(\mathbf{M} \otimes_{\mathbf{A}} \mathbf{N})\) for two finitely generated projective A-modules M, N; \(\kappa(\mathbf{A})\) is the unit element of \(K(\mathbf{A})\) .
\item
  For every finitely generated projective A-module M and every integer \(p \geqslant 0\) let \(\lambda^p\) denote the element \(\kappa(\bigwedge^p M)\) of \(K(A)\) . Show that
\end{enumerate}

\[
\lambda^ {p} (\mathbf {M} \oplus \mathbf {N}) = \sum_ {r + s = p} \lambda^ {r} (\mathbf {M}). \lambda^ {s} (\mathbf {N}).
\]

Let \(\lambda_{\mathrm{T}}(\mathbf{M})\) denote the element \(\sum_{p=0}^{\infty} \lambda^{p}(\mathbf{M}) \mathrm{T}^{p}\) in the ring \(K(\mathrm{A})[[\mathrm{T}]]\) of formal power series in one indeterminate over the ring \(K(\mathrm{A})\) ; show that

\[
\lambda_ {\mathbf {T}} (\mathbf {M} \oplus \mathbf {N}) = \lambda_ {\mathbf {T}} (\mathbf {M}) \lambda_ {\mathbf {T}} (\mathbf {N});
\]

derive a homomorphism also written \(x \mapsto \lambda_{\mathrm{T}}(x)\) of \(K(\mathrm{A})\) to the multiplicative group of formal power series of \(K(\mathrm{A})[[\mathrm{T}]]\) whose constant term is equal to 1. Let \(\lambda^{p}(x)\) denote the coefficient of \(T^{p}\) in \(\lambda_{\mathrm{T}}(x)\) ; then

\[
\lambda^ {p} (x + y) = \sum_ {r + s = p} \lambda^ {r} (x) \lambda^ {s} (y)
\]

for all \(x, y\) in \(\mathbf{K}(\mathbf{A})\) .

\begin{enumerate}
\def\labelenumi{(\alph{enumi})}
\setcounter{enumi}{2}
\item
  If \(\mathbf{A}\) is a local ring, the ring \(K(\mathbf{A})\) is canonically identified with \(\mathbf{Z}\) and \(\lambda^p(n) = \binom{n}{p}\). Deduce that, if \(\mathbf{A}\) is the direct composition of \(m\) local rings, \(K(\mathbf{A})\) is isomorphic to \(\mathbf{Z}^m\) .
\item
  Let \(\phi: \mathbf{A} \to \mathbf{B}\) be a ring homomorphism. Show that there exists a unique ring homomorphism \(\phi': \mathbf{K}(\mathbf{A}) \to \mathbf{K}(\mathbf{B})\) such that \(\phi'(\kappa_{\mathbf{A}}(\mathbf{M})) = \kappa_{\mathbf{B}}(\mathbf{M}_{(\mathbf{B})})\) for every finitely generated projective A-module \(\mathbf{M}\) ; then \(\phi'(\lambda^p(x)) = \lambda^p(\phi'(x))\) for all \(x \in \mathbf{K}(\mathbf{A})\) ; if \(\psi: \mathbf{B} \to \mathbf{C}\) is another ring homomorphism, \((\psi \circ \phi)' = \psi' \circ \phi'\) .
\item
  Let \(\phi: A \to B\) be a ring homomorphism for which \(B\) is a finitely generated projective \(A\) -module; for every finitely generated projective \(B\) -module \(N\) , \(\phi_{*}(N)\) is then a finitely generated projective \(A\) -module; deduce that there exists a \(Z\) -module homomorphism \(\phi_!: K(B) \to K(A)\) such that
\end{enumerate}

\[
\phi_! (\kappa_ {\mathrm{B}} (\mathrm{N})) = \kappa_ {\mathrm{A}} (\phi_ {*} (\mathrm{N})).
\]

Show that, for all \(x \in K(A)\) and \(y \in K(B)\) , \(\phi_!(y \cdot \phi'(x)) = \phi_!(y) \cdot x\) . If \(\psi: B \to C\) is another ring homomorphism such that \(C\) is a finitely generated projective \(B\) -module, then \((\psi \circ \phi)_! = \phi_! \circ \psi_!\) .

\begin{enumerate}
\def\labelenumi{\arabic{enumi}.}
\setcounter{enumi}{10}
\tightlist
\item
  Let K be a field of characteristic \(\neq2\) ; consider a polynomial \(g(\mathbf{X})\in\mathbf{K}[\mathbf{X}]\) of even degree \(\geqslant2\) with all its roots distinct (in an algebraically closed extension of K) and satisfying \(g(0)\neq0\) . Set B = K{[}X{]},
\end{enumerate}

\[
\mathbf {A} = \mathbf {B} [ \mathbf {Y} ] / (\mathbf {Y} ^ {2} - f (\mathbf {X}))
\]

where \(f(\mathbf{X}) = \mathbf{X}g(\mathbf{X})\) .

\begin{enumerate}
\def\labelenumi{(\alph{enumi})}
\item
  Let \(\pmb{y}\) be the class of \(\mathbf{Y}\) in \(\mathbf{A}\) . Show that every \(a \in \mathbf{A}\) may be written uniquely in the form \(a = \mathbf{P} + y\mathbf{Q}\) , where \(\mathbf{P}\) and \(\mathbf{Q}\) are in \(\mathbf{B}\) . Let \(\bar{a} = \mathbf{P} - y\mathbf{Q}\) and \(\mathrm{N}(a)=a\bar{a}=\mathrm{P}^{2}-f\mathrm{Q}^{2}\) . Show that the relation \(\mathrm{N}(a)=0\) implies a=0; deduce that A is an integral domain. Show that the invertible elements of A are the elements \(\neq0\) of K.
\item
  Let \(m\) be the ideal \(AX + Ay\) of \(A\) ; show that \(m\) is maximal and that the ideal \(m_m\) of \(A_m\) is generated by \(y\) (note that \(X = y^2 / g(X)\) ). Deduce that \(m\) is an invertible \(A\) -module (use Theorem 1). Show that \(m^2 = AX\) ; deduce that \(m\) is not a principal ideal of \(A\) and that the group of classes of invertible \(A\) -modules in the field of fractions of \(A\) is not reduced to the identity element. (If \(m = At\) , then \(X = \lambda t^2\) where \(\lambda \in K\) , whence \(\lambda^{-1}X = P^2 + fQ^2\) where \(P \in B\) and \(Q \in B\) ; prove that such a relation is impossible.)
\end{enumerate}

¶ 12. Let A be a ring, I an A-module and B the unique ring whose underlying additive group is A ⊕ I, A being a subring of B and I an ideal of zero square, whose A-module structure obtained by restriction of the scalars is the given structure (Algebra, Chapter II, § 1, Exercise 7).

\begin{enumerate}
\def\labelenumi{(\alph{enumi})}
\item
  Suppose that every non-invertible element of A belongs to the annihilator of some element \(\neq0\) of I. Show that every element in B which is not a divisor of 0 is invertible, in other words that B is equal to its total ring of fractions.
\item
  Let \((\mathfrak{m}_{\lambda})_{\lambda \in L}\) be a family of maximal ideals of A such that every non-invertible ideal of A belongs to at least one \(\mathfrak{m}_{\lambda}\) . Show that, if I is taken to be the direct sum of the A-modules \(\mathrm{A} / \mathfrak{m}_{\lambda}\) , condition (a) is satisfied.
\item
  Deduce from (b) that, if there exists a non-free projective A-module of rank 1, I can be chosen such that B is equal to its total ring of fractions but there exists a non-free projective B-module of rank 1 (this module is necessarily not invertible since B is the only non-degenerate sub-B-module of B).
\item
  Suppose that there exists in A a maximal ideal m such that, for all \(x \in m\) , there exists a maximal ideal \(m' \neq m\) of A containing x. If I is taken to be the A-module the direct sum of the A/m', where \(m'\) runs through the set of maximal ideals of A distinct from m, the corresponding ring \(B = A \oplus I\) is equal to its total ring of fractions; if we write \(n = m \oplus I\) , which is a maximal ideal of B, show that \(B_n\) is isomorphic to \(A_m\) ; if A is an integral domain (necessarily not a field by hypothesis), \(B_n\) is then an integral domain distinct from its field of fractions. If further \(mA_m\) is a principal ideal of the integral domain \(A_m\) , show that the ideal \(mB = m \otimes_A B\) of B is a non-free degenerate projective B-module of rank 1. (We shall see in Chapter VII that there are Dedekind domains A in which there exist maximal ideals m none of whose powers is a principal ideal; for such a ring A, all the preceding hypotheses are satisfied.)
\end{enumerate}

¶ 13. Let A be a (commutative) ring and B an A-algebra (not necessarily commutative). Suppose that the A-module B is faithful and generated by a finite number of elements \(b_i\) ( \(1 \leqslant i \leqslant n\) ); recall that \(B^0\) denotes the opposite algebra of B and that a canonical homomorphism is defined of the A-algebra

B \(\otimes_{\mathbf{A}}\mathbf{B}^{0}\) to the A-algebra \(\operatorname{End}_{\mathbf{A}}(\mathbf{B})\) mapping \(b\otimes b'\) to the endomorphism \(z\mapsto bzb'\) .

\begin{enumerate}
\def\labelenumi{(\alph{enumi})}
\tightlist
\item
  Show that the two following conditions are equivalent:
\end{enumerate}

\((\alpha)\) The \(b_{i}\) form a basis of the A-module B and the canonical homomorphism \(\mathbf{B} \otimes_{\mathbf{A}} \mathbf{B}^{0} \to \operatorname{End}_{\mathbf{A}}(\mathbf{B})\) is bijective.

(β) The matrix \(U = (u_{ij})\) such that \(u_{ij} = b_j b_i\) is invertible in \(\mathbf{M}_n(\mathbf{B})\) .

\begin{enumerate}
\def\labelenumi{(\alph{enumi})}
\setcounter{enumi}{1}
\tightlist
\item
  Suppose further that A is a local ring with maximal ideal m; let \(\beta_{i}\) denote the canonical image of \(b_{i}\) in B/mB. Show that conditions ( \(\alpha\) ) and ( \(\beta\) ) are then equivalent to each of the conditions:
\end{enumerate}

(γ) The \(\beta_{i}\) form a basis of the A/m-module B/mB and B/mB is a simple central A/m-algebra (Algebra, Chapter VIII, § 5, no. 4).

(δ) The matrix \((\beta_{j}\beta_{i})\) is invertible in \(\mathbf{M}_{n}(\mathrm{B}/\mathrm{mB})\) .

(To prove that (δ) implies (β), use the fact that mB is contained in the Jacobson radical of B and Exercise 5 of Algebra, Chapter VIII, § 6).

\begin{enumerate}
\def\labelenumi{\arabic{enumi}.}
\setcounter{enumi}{13}
\tightlist
\item
  Let A be a (commutative) ring and B an A-algebra such that the A-module B is faithfully flat and finitely presented. Show that the following conditions are equivalent:
\end{enumerate}

(α) For every maximal ideal m of A, the A/m-algebra B/mB is simple and central.

(β) B is a projective A-module and, if \(B^{0}\) denotes the opposite algebra to B, the canonical homomorphism \(B \otimes_{A} B^{0} \to \operatorname{End}_{A}(B)\) is bijective.

(Reduce it to the case where A is a local ring and apply Exercise 13.) B is then called an Azumaya algebra over A.

¶ 15. Let B be an Azumaya A-algebra (Exercise 14).

\begin{enumerate}
\def\labelenumi{(\alph{enumi})}
\item
  Show that the centre of B is identical with the subring A of B and is a direct factor of the A-module B (cf.~Chapter I, § 2, Exercise 21 and Chapter II, § 5, Exercise 4).
\item
  Show that, for every two-sided ideal \(\mathfrak{b}\) of B, \(\mathfrak{b} = (\mathfrak{b} \cap A)B\) . (Note first that B is stable under every homomorphism of the A-module B and use the existence of a projector of B onto A.)
\end{enumerate}

\begin{enumerate}
\def\labelenumi{\arabic{enumi}.}
\setcounter{enumi}{15}
\tightlist
\item
  \begin{enumerate}
  \def\labelenumii{(\alph{enumii})}
  \tightlist
  \item
    Let A be a (commutative) ring, A' a commutative A-algebra which is a flat A-module and B an A-algebra. Show that, if B is an Azumaya algebra, \(B' = A' \otimes_{A} B\) is an Azumaya A'-algebra; the converse is true if \(A'\) is assumed to be a faithfully flat A-module (cf.~Algebra, Chapter II, §5, no. 1, Corollary to Proposition 4 and Commutative Algebra, Chapter I, §2, no. 9, Proposition 10 and §3, no. 6, Proposition 12).
  \end{enumerate}
\end{enumerate}

\begin{enumerate}
\def\labelenumi{(\alph{enumi})}
\setcounter{enumi}{1}
\tightlist
\item
  Show that the tensor product \(\mathbf{B} \otimes_{\mathbf{A}} \mathbf{C}\) of two Azumaya A-algebras \(\mathbf{B}, \mathbf{C}\) is an Azumaya A-algebra (cf.~Algebra, Chapter VIII, § 7, no. 4, Corollary 2 to Theorem 2).
\end{enumerate}

¶ 17. Let A be a ring and P a faithful finitely generated projective A-module.

\begin{enumerate}
\def\labelenumi{(\alph{enumi})}
\item
  Show that the A-algebra \(\mathbf{B} = \operatorname{End}_{\mathbf{A}}(\mathbf{P})\) is an Azumaya algebra. (Observe first that the A-module B is isomorphic to \(\mathbf{P} \otimes_{\mathbf{A}} \mathbf{P}^*\) and hence finitely presented, then use §2, no. 7, Proposition 19.)
\item
  Show that P is a left projective B-module (use Theorem 1 (d) and §3, Exercise 8 (c)).
\item
  Let \(\mathbf{P}'\) be another faithful finitely generated projective A-module. Show that \(\mathbf{P} \otimes_{\mathbf{A}} \mathbf{P}'\) is a faithful finitely generated projective A-module and that \(\operatorname{End}_{\mathbf{A}}(\mathbf{P} \otimes_{\mathbf{A}} \mathbf{P}')\) is canonically identified with \(\operatorname{End}_{\mathbf{A}}(\mathbf{P}) \otimes_{\mathbf{A}} \operatorname{End}_{\mathbf{A}}(\mathbf{P}')\) (use Theorem 1 (e), § 2, no. 7, Propositions 18 and 19 and § 3, no. 3, Theorem 1). Consider the case where \(\mathbf{P}'\) is of rank 1.
\end{enumerate}

¶ 18. Let A be a ring, P a faithful finitely generated projective A-module, B the A-algebra \(\operatorname{End}_{\mathbf{A}}(\mathbf{P})\) , E a left B-module and F the A-module \(\operatorname{Hom}_{\mathbf{B}}(\mathbf{P}, \mathbf{E})\) . Show that the canonical homomorphism

\[
\beta \colon \mathrm{P} \otimes_ {\mathbf {A}} \mathrm{F} \rightarrow \mathrm{E}
\]

defined in Algebra, Chapter VIII, § 1, no. 4, is bijective. (Reduce it to the case where A is a local ring as in Exercise 17 (c); in this case apply Exercise 9 of Algebra, Chapter VIII, § 1, observing that \(\mathbf{F}^n\) and \(\operatorname{Hom}_{\mathbf{B}}(\mathbf{P}, \mathbf{E}^n)\) are canonically isomorphic A-modules.)

Show that there exists a strictly increasing bijection of the set of sub-A-modules of F onto the set of sub-B-modules of E (same method).

¶ 19. Let A be a ring, B and C two A-algebra and \(\phi: B \to C\) an A-algebra homomorphism. Suppose that B is an Azumaya algebra (Exercise 14). Let \(C'\) be the sub-A-algebra of C consisting of the elements which commute with the elements of \(\phi(B)\) . Show that C is isomorphic to the algebra \(B \otimes_A C'\) . (Consider C as a \(B \otimes_A B^0\) -module and apply Exercise 18, taking account of the fact that \(B \otimes_A B^0\) is isomorphic to \(\operatorname{End}_A(B)\) .)

¶ 20. Let A be a ring and P, P' two faithful finitely generated projective A-modules. Show that, if θ: End\_A(P) → End\_A(P') is an A-algebra isomorphism, there exists a projective A-module F of rank 1 and an A-isomorphism φ: P' → P ⊗\_A F such that \(\theta(u) = \phi^{-1} \circ (u \otimes 1) \circ \phi\) for all u ∈ End\_A(P). (Apply Exercise 18 to E = P' considered as an (End\_A(P))-module by means of θ; to prove that Hom\_B(P, P') is a projective A-module of rank 1, use Theorem 2 of no. 3 and Algebra, Chapter VIII, § 5, Exercise 7 (c).)

¶ 21. Let A be a ring and n an integer ≥0.

\begin{enumerate}
\def\labelenumi{(\alph{enumi})}
\tightlist
\item
  Show that A-modules F such that \(F^{n}\) is isomorphic to \(A^{n}\) are projective of rank 1 and that their classes in \(P(A)\) form a subgroup \(P_{n}(A)\) every element of which is annihilated by n.~(To show the first assertion, use Proposition
\end{enumerate}

5 of §2, no. 2; for the second, observe that \(\bigwedge^{n} \mathbf{F}^{n}\) and \(\otimes \mathbf{F}\) are canonically isomorphic.)

\begin{enumerate}
\def\labelenumi{(\alph{enumi})}
\setcounter{enumi}{1}
\tightlist
\item
  Let \(\operatorname{Aut}_n(A)\) be the group of automorphisms of the A-algebra of matrices \(\mathbf{M}_n(A)\) and let \(\operatorname{Int}_n(A)\) be the subgroup of inner automorphisms of \(\mathbf{M}_n(A)\) . Show that the quotient group \(\operatorname{Aut}_n(A)/\operatorname{Int}_n(A)\) is isomorphic to \(P_n(A)\) (use Exercise 20 with \(P = P' = A^n\) or Exercise 9 of Algebra, Chapter VIII, § 1). Deduce that, if \(A\) is a semi-local ring or a principal ideal domain, then
\end{enumerate}

\[
\operatorname{Aut} _ {n} (\mathbf {A}) = \operatorname{Int} _ {n} (\mathbf {A}).
\]

* (c) If A is a Dedekind domain, show that \(P_{n}(A)\) is identical with the subgroup of \(P(A)\) consisting of the elements annihilated by \(n_*\)

¶ 22. Let A be a ring and M a finitely generated projective A-module. For a finitely presented submodule M' of M to be a direct factor of M, it is necessary and sufficient that it be a pure submodule of M (use Theorem 1 and § 3, Exercise 15).

\begin{enumerate}
\def\labelenumi{\arabic{enumi}.}
\setcounter{enumi}{22}
\tightlist
\item
  Let A be a ring and \((f_i)_{i\in I}\) a finite family of elements of A generating the ideal A.
\end{enumerate}

\begin{enumerate}
\def\labelenumi{(\alph{enumi})}
\item
  Let \(\mathbf{M}\) be an A-module. Show that, if, for all \(i\in \mathbf{I}\) , \(\mathbf{M}_{f_i}\) is a finitely generated \(\mathbf{A}_{f_i}\) -module, \(\mathbf{M}\) is a finitely generated A-module.
\item
  Let \(\mathbf{B}\) be a (commutative) A-algebra. Show that, if, for all \(i\in \mathbf{I}\) , \(\mathbf{B}_{f_i}\) is a finitely generated \(\mathbf{A}_{f_i}\) -algebra, \(\mathbf{B}\) is a finitely generated A-algebra.
\end{enumerate}

\begin{enumerate}
\def\labelenumi{\arabic{enumi}.}
\setcounter{enumi}{23}
\tightlist
\item
  Let \(\mathbf{A}\) be a ring such that \(\operatorname{Spec}(\mathbf{A})\) is connected.
\end{enumerate}

\begin{enumerate}
\def\labelenumi{(\alph{enumi})}
\item
  If M is a projective A-module of rank n, every direct factor N of M not reduced to 0 admits a rank \(\leqslant n\) .
\item
  Let \(u\) be a homomorphism of a projective A-module of finite rank \(M\) to a projective A-module of finite rank \(M'\) such that \(u(M)\) is a direct factor of \(M'\) . Show that \(\operatorname{Ker}(u)\) is then a direct factor of \(M\) (cf.~Exercise 23) and
\end{enumerate}

\[
\operatorname{rg} (\operatorname{Ker} (u)) + \operatorname{rg} (\operatorname{Im} (u)) = \operatorname{rg} (\mathbf {M}).
\]

Moreover the transpose \({}^t u\) is such that \({}^t u(\mathbf{M}^*)\) is a direct factor of \(\mathbf{M}^*\) and \(\operatorname{rg}({}^t u(\mathbf{M}^*)) = \operatorname{rg}(u(\mathbf{M}))\) (cf.~§ 3, Exercise 14 (a)).

\begin{enumerate}
\def\labelenumi{\arabic{enumi}.}
\setcounter{enumi}{24}
\tightlist
\item
  Let A be a ring and S a multiplicative subset of A. An ideal a of A is called S-invertible if there exists \(s \in S\) and an ideal b of A such that ab = As; two ideals a, b of A are called S-equivalent if there exists s, t in S such that sa = tb. Show that the S-equivalence classes of the S-invertible ideals of A form a multiplicative group \(C_{S}\) . If S is saturated ( \(\S 2\) , Exercise 1) and no element of S is a divisor of 0 in A, show that the group \(C_{S}\) is isomorphic to the classes of invertible sub- \(S^{-1}A\) -modules of the total ring of fractions B of A (and \(S^{-1}A\) ).
\end{enumerate}

CHAPTER III(*)

